\documentclass{article}
\usepackage[final]{neurips_2024}

\usepackage[utf8]{inputenc}
\usepackage[T1]{fontenc}
\usepackage{hyperref}
\usepackage{url}
\usepackage{booktabs}
\usepackage{amsfonts}
\usepackage{amsmath}
\usepackage{amssymb}
\usepackage{nicefrac}
\usepackage{microtype}
\usepackage{graphicx}
\usepackage{xcolor}
\usepackage{multirow}
\usepackage{adjustbox}

\newcommand{\ours}{\textsc{C2UD}}
\newcommand{\RS}{\mathrm{RS}}
\newcommand{\CD}{\mathrm{CD}}
\newcommand{\PA}{\mathrm{PA}}
\newcommand{\JSD}{\mathrm{JSD}}
\newcommand{\NLI}{\mathrm{NLI}}

\title{Context-Contrastive Uncertainty Decomposition for Reliable Retrieval-Augmented Generation}

\author{
  Anonymous Author(s)
}

\begin{document}
\maketitle

\begin{abstract}
Retrieval-Augmented Generation (RAG) has become the dominant paradigm for grounding large language models in external knowledge, yet estimating whether a RAG system's output is trustworthy remains an open problem.
Existing uncertainty estimation methods produce a single score that conflates fundamentally different failure modes---retrieval failure, grounding failure, and parametric override.
We propose \ours{} (Context-Contrastive Uncertainty Decomposition), which decomposes RAG uncertainty into three interpretable components by comparing an LLM's behavior under three controlled conditions: with retrieved context, without context, and with irrelevant context.
This three-condition design enables a \emph{Context Discrimination} metric that detects indiscriminate context sensitivity---a failure mode invisible to two-condition methods.
We evaluate \ours{} on Natural Questions, TriviaQA, and PopQA across three model families (Llama~3.1~8B, Mistral~7B, Phi-3~Mini).
Averaged across models, \ours{} achieves AUROC of 0.668 on TriviaQA and 0.689 on PopQA for selective prediction, outperforming the strongest context-aware baseline CRUX (0.580 and 0.641 respectively).
Performance varies substantially across models: on Mistral, \ours{} achieves 0.797 AUROC on TriviaQA, while on Llama, simpler baselines remain competitive.
We provide honest analysis of when the three-condition framework helps and when it does not, along with a failure mode diagnosis that achieves significant component separation on Mistral.
\end{abstract}

\section{Introduction}

Retrieval-Augmented Generation (RAG) has become the standard approach for grounding large language model (LLM) outputs in external knowledge, enabling more factual and current responses~\citep{lewis2020rag}.
However, RAG systems fail in multiple distinct ways: the retriever may return irrelevant documents, the generator may hallucinate despite good context, or the model may override retrieved evidence with incorrect parametric knowledge.
A critical open problem is how to reliably estimate whether a RAG system's output is trustworthy.

Recent axiomatic analysis by \citet{soudani2025why} has formally demonstrated that existing uncertainty estimation (UE) methods---including token probability, semantic entropy~\citep{farquhar2024detecting}, and verbalized confidence---systematically fail in RAG settings.
These methods produce a single, monolithic uncertainty score that conflates fundamentally different failure modes.
While concurrent work has begun addressing context-awareness in uncertainty estimation---notably CRUX~\citep{yuan2026crux}, which compares model behavior with and without context---these approaches use only \emph{two} conditions (with vs.\ without context) and cannot distinguish between a model that genuinely grounds in relevant evidence versus one that is indiscriminately sensitive to any context.

Our key insight is that by comparing an LLM's behavior under \textbf{three} controlled conditions---with retrieved context, without context, and with irrelevant context---we can disentangle failure modes that two-condition approaches conflate.
This is analogous to a controlled experiment: the ``with context'' setting is the treatment, ``without context'' is the control, and ``irrelevant context'' is the \emph{placebo}.
The pattern of divergence across these three conditions reveals not just \emph{whether} the model uses context, but whether it uses \emph{the right} context.

We propose \ours{} (Context-Contrastive Uncertainty Decomposition), which decomposes RAG uncertainty into three interpretable components:
\textbf{Retrieval Sensitivity} (RS), measuring how much the model's output changes when context is provided;
\textbf{Context Discrimination} (CD), measuring whether the model differentially responds to relevant versus irrelevant context; and
\textbf{Parametric Agreement} (PA), measuring semantic agreement between context-grounded and parametric-only answers.

Our contributions are:
\begin{itemize}
    \item We propose \ours{}, a three-condition contrastive framework that decomposes RAG uncertainty into three interpretable components, enabling the novel Context Discrimination metric for detecting indiscriminate context sensitivity.
    \item We evaluate \ours{} on three open-domain QA benchmarks across three model families and seven baselines, showing that the three-condition design provides significant improvements on TriviaQA and PopQA (averaged across models), with particularly strong results on Mistral (AUROC 0.797 on TriviaQA).
    \item We provide transparent analysis of model-dependent performance variation and demonstrate that \ours{}'s failure mode diagnosis achieves statistically significant component separation on Mistral ($p < 10^{-10}$ for CD between failure types), while separation is weaker on other models.
\end{itemize}

The remainder of this paper is organized as follows: Section~\ref{sec:related} reviews related work, Section~\ref{sec:method} describes the \ours{} framework, Section~\ref{sec:experiments} presents our experimental evaluation, Section~\ref{sec:discussion} discusses limitations, and Section~\ref{sec:conclusion} concludes.

\section{Related Work}
\label{sec:related}

\paragraph{Uncertainty Estimation for LLMs.}
Semantic Entropy~\citep{kuhn2023semantic,farquhar2024detecting} clusters sampled responses by semantic equivalence and computes entropy over clusters, separating linguistic from semantic uncertainty.
Semantic Entropy Probes~\citep{kossen2024semantic} approximate this from hidden states in a single pass but suffer from out-of-distribution degradation.
SPUQ~\citep{gao2024spuq} perturbs input prompts to estimate epistemic uncertainty.
Self-Consistency~\citep{wang2023selfconsistency} uses majority vote across sampled reasoning paths.
\citet{kadavath2022language} study LLMs' ability to evaluate their own knowledge.
All these methods operate on the generation level and do not account for RAG-specific interactions between context and generation.

\paragraph{Uncertainty in RAG.}
\citet{soudani2025why} present an axiomatic analysis showing that existing UE methods fail in RAG settings and propose a calibration function.
However, their calibration is a post-hoc fix that does not decompose uncertainty by source.
Counterfactual prompting~\citep{xia2024controlling} uses counterfactual variations to assess RAG risk but focuses on risk control rather than uncertainty decomposition.

\paragraph{Context-Aware Contrastive Methods.}
CRUX~\citep{yuan2026crux} is the most closely related work, proposing a context-aware dual-metric framework using contextual entropy reduction and unified consistency examination with two conditions (with/without context).
\ours{} differs in three key ways: (1)~\ours{} uses three conditions (adding irrelevant context), enabling the Context Discrimination metric; (2)~\ours{} uses text-level comparison rather than sampling-based entropy, requiring fewer forward passes; and (3)~\ours{} explicitly decomposes uncertainty into three named components with distinct diagnostic interpretations.
FRANQ~\citep{fadeeva2025franq} decomposes RAG uncertainty conditioned on faithfulness to retrieved context, operating post-hoc at the claim level, while \ours{} operates at the generation-process level by controlling input conditions before generation.
Context-Aware Decoding~\citep{shi-etal-2024-trusting} contrasts output distributions with and without context to improve faithfulness; \ours{} extends this contrastive principle from decoding to uncertainty estimation with a third control condition.

\section{Method}
\label{sec:method}

\subsection{Overview}

Given a query $q$ and retrieved documents $D = \{d_1, \ldots, d_k\}$, \ours{} generates responses under three controlled conditions, extracts three uncertainty components, and combines them via lightweight calibration.
The decomposition itself is training-free and model-agnostic; only the final scoring aggregation requires a small labeled calibration set (${\sim}250$ examples).

\subsection{Multi-Condition Generation}

We generate responses under three conditions:
\begin{enumerate}
    \item \textbf{RAG condition}: Present the query with retrieved documents. Obtain response $a_D$.
    \item \textbf{Parametric condition}: Present the query alone (no documents). Obtain response $a_\emptyset$.
    \item \textbf{Control condition}: Present the query with irrelevant documents $D'$ sampled from unrelated Wikipedia passages. Obtain response $a_{D'}$.
\end{enumerate}

\subsection{Uncertainty Component Extraction}

\paragraph{Retrieval Sensitivity (RS).}
Measures how much the model's output changes when retrieved context is provided versus absent:
\begin{equation}
    \RS(q, D) = 1 - \mathrm{F1}(a_D, a_\emptyset)
\end{equation}
where $\mathrm{F1}$ is the token-level F1 score between the two responses.
High RS indicates the model's answer changed substantially with context; low RS indicates the context was largely ignored.

\paragraph{Design rationale: text-level vs.\ distributional comparison.}
Our original formulation specified sequence-level Jensen-Shannon Divergence (JSD) over token-level probability distributions for RS and CD computation.
In practice, the top-$k$ logprob distributions available from standard inference APIs are too sparse to produce reliable JSD estimates---with only 20 logprobs per position out of vocabularies of 32k--128k tokens, the ``other'' probability mass dominates and compresses JSD values into a narrow range (std $\approx 0.03$), limiting discriminative power.
We therefore adopt text-level proxies: token F1 for RS and a normalized text-similarity ratio for CD, which produce substantially higher variance (std $\approx 0.13$) and better discriminate between correct and incorrect answers.
We discuss this trade-off further in Section~\ref{sec:discussion}.

\paragraph{Context Discrimination (CD).}
Measures whether the model differentially responds to relevant versus irrelevant context:
\begin{equation}
    \CD(q, D) = \frac{1 - \mathrm{F1}(a_D, a_{D'})}{(1 - \mathrm{F1}(a_D, a_{D'})) + (1 - \mathrm{F1}(a_\emptyset, a_{D'}))}
\end{equation}
This ratio formulation ensures $\CD \in [0, 1]$.
When $\CD > 0.5$, the model's response to relevant documents diverges from its response to irrelevant documents more than the parametric-only response does, suggesting genuine grounding.
When $\CD \leq 0.5$, the model is no more discriminating of relevant context than the baseline, suggesting unreliable context use.
Edge cases where both terms are zero (model ignores all context) default to $\CD = 0.5$.

\paragraph{Parametric Agreement (PA).}
Measures semantic agreement between the RAG answer and the parametric answer using a lightweight NLI model (DeBERTa-v3-base fine-tuned on MNLI/FEVER/ANLI):
\begin{equation}
    \PA(q, D) = P(\text{entailment} \mid a_D, a_\emptyset) - P(\text{contradiction} \mid a_D, a_\emptyset)
\end{equation}
We use the entailment-minus-contradiction margin rather than raw entailment probability.
The raw entailment probability is heavily right-skewed for short QA answers (most pairs receive high entailment scores), compressing the useful signal.
The margin produces a less skewed distribution with better discriminative power.

\subsection{Calibrated Uncertainty Scoring}

We feed the three components plus pairwise interactions ($\RS \times \CD$, $\RS \times \PA$, $\CD \times \PA$) as six features to a logistic regression classifier for selective prediction.
This classifier is fitted on a small labeled calibration set (${\sim}250$ examples with binary correctness labels), separate from the test set, using 5-fold cross-validation to select regularization strength $C \in \{0.01, 0.1, 1, 10, 100\}$.
This is analogous to how Platt scaling or temperature scaling calibrates any confidence method.

\subsection{Targeted Intervention}

Based on the decomposed uncertainty, \ours{} enables specific remediation:
\begin{itemize}
    \item \textbf{Low RS} (context ignored): The retrieval may be irrelevant---try a different retrieval strategy or more passages.
    \item \textbf{Low CD} (poor grounding): Trigger re-retrieval with a reformulated query.
    \item \textbf{Low PA with high CD} (knowledge conflict): The model is grounded but parametric knowledge conflicts---flag for human review or prioritize the context-grounded answer.
\end{itemize}

\section{Experiments}
\label{sec:experiments}

\subsection{Setup}

\paragraph{Datasets.}
We evaluate on three standard open-domain QA benchmarks: \textbf{Natural Questions (NQ)}~\citep{kwiatkowski2019natural} (500 test examples sampled from the open-domain subset), \textbf{TriviaQA}~\citep{joshi2017triviaqa} (500 examples), and \textbf{PopQA}~\citep{mallen2023when} (500 examples, entity-centric long-tail questions where RAG is particularly important).
For retrieval, we use BM25 over a Wikipedia corpus, retrieving top-5 passages per query.
Irrelevant control documents are sampled from unrelated Wikipedia topic categories with no keyword overlap with the query.

\paragraph{Models.}
We test across three model families: \textbf{Llama~3.1~8B Instruct} (Meta), \textbf{Mistral~7B Instruct v0.3} (Mistral AI), and \textbf{Phi-3~Mini~3.8B Instruct} (Microsoft).
All models use greedy decoding (temperature 0) via vLLM~\citep{izacard2022contriever} on a single NVIDIA A6000 (48GB).

\paragraph{Baselines.}
We compare against seven methods:
(1)~\textbf{Token Probability}: mean log-probability of generated tokens;
(2)~\textbf{Verbalized Confidence}: prompted self-reported confidence (0--100\%);
(3)~\textbf{Semantic Entropy}~\citep{farquhar2024detecting}: entropy over semantic clusters from 5 sampled responses;
(4)~\textbf{Self-Consistency}~\citep{wang2023selfconsistency}: largest semantic cluster fraction from 5 samples;
(5)~\textbf{SPUQ}~\citep{gao2024spuq}: perturbation-based uncertainty via prompt paraphrasing;
(6)~\textbf{Axiomatic Calibration}~\citep{soudani2025why}: calibrated token probability for RAG;
(7)~\textbf{CRUX}~\citep{yuan2026crux}: two-condition contextual entropy reduction + consistency, combined via logistic regression with the same calibration procedure as \ours{} for fair comparison.

\paragraph{Metrics.}
We evaluate selective prediction via \textbf{AUROC} (area under ROC curve for predicting answer correctness), \textbf{AUPRC} (area under precision-recall curve), and \textbf{Coverage@90} (maximum fraction of questions answered at $\geq$90\% accuracy).
We report means and standard deviations over 3 random seeds (42, 43, 44) for 50/50 calibration/test splits.
Answer correctness uses token-level F1 with threshold 0.5 (standard SQuAD-style evaluation).

\subsection{Main Results}

\begin{table}[t]
\caption{AUROC ($\uparrow$) for selective prediction, by model and dataset. Best per column in \textbf{bold}. \ours{} (full) uses all three components (RS+CD+PA); results are means over 3 seeds. TQA = TriviaQA, PQA = PopQA.}
\label{tab:main_per_model}
\centering
\begin{adjustbox}{max width=\textwidth}
\begin{tabular}{l ccc c ccc c ccc}
\toprule
& \multicolumn{3}{c}{Llama 3.1 8B} && \multicolumn{3}{c}{Mistral 7B} && \multicolumn{3}{c}{Phi-3 Mini} \\
\cmidrule{2-4} \cmidrule{6-8} \cmidrule{10-12}
Method & NQ & TQA & PQA && NQ & TQA & PQA && NQ & TQA & PQA \\
\midrule
Token Prob. & 0.492 & 0.513 & 0.430 && 0.607 & 0.589 & \textbf{0.722} && 0.573 & 0.499 & 0.636 \\
Verbalized & \textbf{0.676} & \textbf{0.656} & \textbf{0.700} && 0.601 & 0.548 & 0.604 && 0.522 & 0.531 & 0.455 \\
Sem.\ Entropy & 0.503 & 0.537 & 0.459 && 0.496 & 0.447 & 0.443 && 0.483 & 0.552 & 0.504 \\
Self-Consist. & 0.505 & 0.535 & 0.461 && 0.496 & 0.447 & 0.443 && 0.483 & 0.552 & 0.504 \\
Axiomatic & 0.481 & 0.541 & 0.480 && 0.603 & 0.583 & 0.698 && \textbf{0.596} & 0.569 & 0.670 \\
CRUX & 0.588 & 0.554 & 0.693 && 0.503 & 0.528 & 0.669 && 0.484 & 0.651 & 0.556 \\
\midrule
\ours{} (full) & 0.593 & 0.581 & 0.694 && 0.583 & \textbf{0.797} & 0.751 && 0.536 & \textbf{0.632} & \textbf{0.680} \\
\bottomrule
\end{tabular}
\end{adjustbox}
\end{table}

Table~\ref{tab:main_per_model} presents AUROC results for selective prediction across all model-dataset combinations.
The results reveal substantial model-dependent variation:

\textbf{Mistral~7B} shows the strongest \ours{} performance.
On TriviaQA, \ours{} achieves 0.797 AUROC, outperforming all baselines by a wide margin (next best: Token Probability at 0.589, $\Delta = +0.208$).
On PopQA, \ours{} achieves 0.751 versus CRUX's 0.669 ($\Delta = +0.082$).
On NQ, \ours{} (0.583) trails the RS-only variant (0.617) and Token Probability (0.607).

\textbf{Llama~3.1~8B} presents a contrasting picture: Verbalized Confidence dominates across all three datasets (0.676, 0.656, 0.700), outperforming both \ours{} and CRUX.
\ours{} does outperform CRUX on NQ (0.593 vs.\ 0.588) and PopQA (0.694 vs.\ 0.693), but the margins are within standard error.
This indicates that on Llama, the contrastive framework provides limited benefit over well-calibrated verbalized self-assessment.

\textbf{Phi-3~Mini} shows intermediate results: \ours{} is best on TriviaQA (0.632) and PopQA (0.680), but Axiomatic Calibration leads on NQ (0.596).

\begin{table}[t]
\caption{Cross-model averaged results. AUROC and AUPRC ($\uparrow$). Best in \textbf{bold}.}
\label{tab:main_avg}
\centering
\begin{adjustbox}{max width=\textwidth}
\begin{tabular}{l cc cc cc}
\toprule
& \multicolumn{2}{c}{NQ} & \multicolumn{2}{c}{TriviaQA} & \multicolumn{2}{c}{PopQA} \\
\cmidrule(lr){2-3} \cmidrule(lr){4-5} \cmidrule(lr){6-7}
Method & AUROC & AUPRC & AUROC & AUPRC & AUROC & AUPRC \\
\midrule
Token Prob. & 0.557 & \textbf{0.375} & 0.534 & 0.689 & 0.596 & 0.360 \\
Verbalized & \textbf{0.600} & 0.343 & 0.578 & 0.655 & 0.587 & 0.260 \\
Sem.\ Entropy & 0.492 & 0.286 & 0.507 & 0.617 & 0.479 & 0.204 \\
Axiomatic & 0.560 & 0.347 & 0.564 & 0.666 & 0.616 & 0.327 \\
\midrule
CRUX & 0.522 & 0.309 & 0.580 & 0.665 & 0.641 & 0.346 \\
\ours{} (full) & 0.552 & 0.372 & \textbf{0.668} & \textbf{0.750} & \textbf{0.689} & \textbf{0.412} \\
\bottomrule
\end{tabular}
\end{adjustbox}
\end{table}

Table~\ref{tab:main_avg} shows cross-model averaged results.
\ours{} achieves the best AUROC on TriviaQA (0.668 vs.\ CRUX 0.580, $\Delta = +0.088$) and PopQA (0.689 vs.\ CRUX 0.641, $\Delta = +0.048$).
On NQ, Verbalized Confidence leads (0.600), with \ours{} at 0.552.
The TriviaQA advantage is primarily driven by Mistral's strong performance.

\paragraph{Coverage@90.}
Coverage@90 results are near zero for almost all method-model-dataset combinations.
The highest value is \ours{} on Mistral/TriviaQA (Cov@90 = 0.308), reflecting that base QA accuracy ranges from 20--60\% depending on model and dataset, making 90\% selective accuracy nearly unattainable at meaningful coverage.

\subsection{Ablation Study}

\begin{table}[t]
\caption{Ablation study: AUROC by \ours{} variant, averaged across all three models. RS+PA is the two-condition variant analogous to CRUX's information set (no irrelevant context condition).}
\label{tab:ablation}
\centering
\begin{tabular}{lcccc}
\toprule
Variant & NQ & TriviaQA & PopQA & Avg \\
\midrule
RS only & 0.588 & 0.656 & 0.686 & 0.643 \\
CD only & 0.507 & 0.612 & 0.553 & 0.557 \\
PA only & 0.477 & 0.544 & 0.618 & 0.547 \\
\midrule
RS+CD & \textbf{0.588} & \textbf{0.682} & 0.676 & \textbf{0.649} \\
RS+PA (two-cond.) & 0.548 & 0.664 & \textbf{0.694} & 0.635 \\
CD+PA & 0.501 & 0.621 & 0.623 & 0.581 \\
Full (RS+CD+PA) & 0.552 & 0.668 & 0.689 & 0.637 \\
\bottomrule
\end{tabular}
\end{table}

Table~\ref{tab:ablation} shows the contribution of each component.
Key findings:

\textbf{RS is the strongest individual component} (avg AUROC 0.643), confirming that the basic contrastive signal---whether context changes the answer---carries the most information for selective prediction.

\textbf{Adding CD to RS improves performance on TriviaQA} ($0.656 \to 0.682$, $\Delta = +0.026$) and achieves the highest average AUROC (0.649), demonstrating that the irrelevant context condition provides useful signal beyond two-condition approaches.

\textbf{The two-condition RS+PA variant} (analogous to CRUX's information set) achieves 0.635 average AUROC.
The full three-condition model (0.637) marginally improves on this, while RS+CD (0.649) provides a larger gain.
This indicates the primary benefit of the third condition flows through the CD component rather than through enriching PA.

\textbf{PA alone is the weakest component} (avg 0.547), likely because the NLI-based comparison is noisy for short QA answers.
However, PA adds value in combination: RS+PA outperforms RS alone on PopQA (0.694 vs.\ 0.686).

\subsection{Failure Mode Analysis}
\label{sec:failure}

We categorize incorrect predictions by failure type using heuristic labels: \emph{retrieval failure} (no gold answer substring in retrieved passages), \emph{grounding failure} (gold answer present but model gives wrong answer), and \emph{parametric override} (parametric answer correct but RAG answer wrong).

\begin{table}[t]
\caption{Mean \ours{} component values by failure type. Statistically significant differences (Welch's $t$-test) marked: $^{**}p < 0.01$, $^{***}p < 0.001$. Sample sizes in parentheses.}
\label{tab:failure}
\centering
\begin{adjustbox}{max width=\textwidth}
\begin{tabular}{l ccc ccc}
\toprule
& \multicolumn{3}{c}{Llama 3.1 8B} & \multicolumn{3}{c}{Mistral 7B} \\
\cmidrule(lr){2-4} \cmidrule(lr){5-7}
Failure Type & RS & CD & PA & RS & CD & PA \\
\midrule
Retrieval (580 / 670) & 0.731 & 0.470 & 0.875 & 0.642 & 0.456 & 0.724 \\
Grounding (32 / 42) & 0.718 & 0.489 & 0.904 & 0.644 & 0.515 & 0.789 \\
Param.\ override (231 / 299) & 0.744 & 0.488$^{**}$ & 0.905 & 0.650 & 0.578$^{***}$ & 0.608$^{***}$ \\
\bottomrule
\end{tabular}
\end{adjustbox}
\end{table}

Table~\ref{tab:failure} shows mean component values by failure type.
On \textbf{Mistral}, CD significantly differentiates retrieval failure from parametric override ($p = 2.95 \times 10^{-11}$): parametric override exhibits higher CD (0.578 vs.\ 0.456), consistent with the model differentially engaging with relevant context when it has competing parametric knowledge.
PA also significantly differentiates these categories ($p = 1.95 \times 10^{-4}$): parametric override shows lower PA (0.608 vs.\ 0.724), as expected when RAG and parametric answers disagree.
Additionally, PA differentiates grounding failure from parametric override ($p = 0.005$).

On \textbf{Llama}, component distributions are more compressed (RS std $\approx 0.12$ vs.\ Mistral's $\approx 0.20$), with only CD between retrieval failure and parametric override reaching significance ($p = 0.002$).
The narrower distributions partially explain both weaker failure mode separation and lower selective prediction performance on this model.

We acknowledge that the grounding failure category is small ($n = 32$--$42$) and that the heuristic labels (based on string matching) introduce noise.
A larger manually-annotated evaluation set would strengthen these conclusions.

\subsection{Targeted Intervention}

\begin{table}[t]
\caption{Intervention experiment on Llama~3.1~8B / NQ (250 test queries). All non-baseline strategies operate at matched coverage (14.4\%).}
\label{tab:intervention}
\centering
\begin{tabular}{lcc}
\toprule
Strategy & Accuracy ($\uparrow$) & Coverage \\
\midrule
No intervention & 0.352 & 1.000 \\
Uniform abstention & 0.583 & 0.144 \\
Uniform re-retrieval & 0.667 & 0.144 \\
\ours{}-intervene & 0.667 & 0.144 \\
\bottomrule
\end{tabular}
\end{table}

Table~\ref{tab:intervention} shows the targeted intervention results.
At matched coverage (14.4\%), \ours{}-intervene and uniform re-retrieval both achieve 66.7\% accuracy, outperforming uniform abstention (58.3\%) by 8.4 percentage points.
The identical performance of targeted and uniform re-retrieval indicates that at this scale and on this model-dataset combination, the decomposition does not provide additional benefit beyond identifying which queries to intervene on.
This experiment is limited (250 queries, single model/dataset) and conducted on Llama/NQ where \ours{}'s decomposition is weakest.
On Mistral, where failure mode separation is much stronger (Table~\ref{tab:failure}), we would expect targeted interventions to show clearer advantages.

\subsection{Calibration Analysis}

Figure~\ref{fig:calibration} shows calibration reliability diagrams for selected methods on Mistral/TriviaQA.
\ours{} achieves ECE of 0.137 and Brier score of 0.204 on Mistral/TriviaQA---competitive with baselines while providing substantially better discrimination (AUROC 0.797).
On Llama/NQ, where component signals are weaker, \ours{}'s ECE is 0.052 with Brier score 0.228.
The logistic regression calibration produces well-calibrated estimates when the underlying components are discriminative but cannot compensate for weak component signal.

\begin{figure}[t]
\centering
\includegraphics[width=0.9\linewidth]{figures/figure7_calibration.png}
\caption{Calibration reliability diagrams for Llama~3.1~8B across three datasets. Diagonal represents perfect calibration. \ours{} (blue) maintains competitive calibration while achieving higher discrimination on TriviaQA and PopQA.}
\label{fig:calibration}
\end{figure}

\section{Discussion and Limitations}
\label{sec:discussion}

\paragraph{Text-Level Proxy vs.\ Distributional Comparison.}
The most significant departure from our original formulation is replacing sequence-level JSD with text-level token F1 for RS and CD.
The proposed JSD formulation, $\JSD_{\text{seq}}(\mathbf{p} \| \mathbf{q}) = \frac{1}{T} \sum_{t=1}^{T} \JSD(p^{(t)} \| q^{(t)})$, requires full next-token distributions at each position.
Standard inference APIs provide only top-$k$ logprobs ($k = 20$), covering a tiny fraction of vocabulary (32k--128k tokens).
The resulting JSD estimates are dominated by the ``other'' mass, compressing values into a narrow range with insufficient variance for discrimination.

The text-based proxy achieves higher variance and better selective prediction, but discards distributional information---two responses with high token F1 may differ substantially in their underlying probability distributions.
This likely contributes to the compressed failure mode distributions: token F1 captures \emph{what} changed but not \emph{how confidently}.
First-token JSD (using only the first token's distribution, where top-$k$ covers more probability mass) is a computationally trivial middle ground we leave to future work.
Access to full vocabulary logprobs (available with some inference frameworks) would enable the original JSD formulation.

\paragraph{Model-Dependent Effectiveness.}
\ours{}'s effectiveness varies substantially by model: strong on Mistral (AUROC 0.797 on TriviaQA), but weaker than Verbalized Confidence on Llama.
This appears related to response variation across conditions: Mistral produces more diverse responses, yielding higher-variance RS and CD components with better discriminative power.
Llama's responses are more consistent across conditions, compressing the contrastive signal.
Practitioners should assess response variation (RS std $> 0.15$ as a heuristic) before adopting \ours{} over simpler alternatives.

\paragraph{PA Formula Deviation.}
Our PA uses entailment-minus-contradiction margin rather than raw entailment probability as originally specified.
The raw probability is heavily right-skewed for short QA answers, compressing useful signal.
The margin produces better discrimination but at the cost of a less interpretable range ($[-1, 1]$ vs.\ $[0, 1]$).

\paragraph{Weak Intervention Evidence.}
The intervention experiment (250 queries, Llama/NQ) is underpowered and conducted on the model-dataset where \ours{} is weakest.
At matched coverage, targeted and uniform re-retrieval achieve identical accuracy, failing to demonstrate the practical value of decomposed interventions.
A more convincing evaluation would use Mistral (where decomposition is strongest) with more queries.

\paragraph{Additional Limitations.}
We evaluate only on open-domain QA; generalization to other RAG tasks (summarization, multi-hop reasoning) is untested.
The 500-example test sets, while standard, limit statistical power.
Heuristic failure mode labels introduce noise.
The original success criterion---significant improvement on at least 2 of 3 datasets---is met for the cross-model average but not for Llama individually, where \ours{} does not consistently outperform Verbalized Confidence.

\section{Conclusion}
\label{sec:conclusion}

We presented \ours{}, a three-condition contrastive framework for decomposing RAG uncertainty into Retrieval Sensitivity, Context Discrimination, and Parametric Agreement.
The key innovation is the irrelevant-context control condition that enables detection of indiscriminate context sensitivity through the Context Discrimination metric.
Averaged across three model families, \ours{} achieves the best selective prediction on TriviaQA (AUROC 0.668, $\Delta = +0.088$ vs.\ CRUX) and PopQA (0.689, $\Delta = +0.048$), with particularly strong results on Mistral (0.797 on TriviaQA).

The framework's effectiveness is model-dependent: models that produce varied responses across context conditions benefit most.
The practical text-based proxy, while effective for selective prediction, limits the granularity of failure mode diagnosis compared to the theoretical distributional formulation.

Future work should (1)~implement first-token JSD or full-vocabulary JSD as distributional alternatives to the text-based proxy, (2)~evaluate on additional RAG tasks and model families, (3)~conduct adequately powered intervention experiments on models where \ours{}'s decomposition is strongest, and (4)~explore adaptive selection of uncertainty estimation methods based on model characteristics.

\bibliographystyle{plainnat}

\begin{thebibliography}{17}

\bibitem[Farquhar et~al.(2024)]{farquhar2024detecting}
Sebastian Farquhar, Jannik Kossen, Lorenz Kuhn, and Yarin Gal.
\newblock Detecting hallucinations in large language models using semantic entropy.
\newblock \emph{Nature}, 630:625--630, 2024.

\bibitem[Fadeeva et~al.(2025)]{fadeeva2025franq}
Ekaterina Fadeeva, Aleksandr Rubashevskii, Dzianis Piatrashyn, Roman Vashurin, Shehzaad Dhuliawala, Artem Shelmanov, Timothy Baldwin, Preslav Nakov, Mrinmaya Sachan, and Maxim Panov.
\newblock Faithfulness-aware uncertainty quantification for fact-checking the output of retrieval augmented generation.
\newblock \emph{arXiv preprint arXiv:2505.21072}, 2025.

\bibitem[Gao et~al.(2024)]{gao2024spuq}
Xiang Gao, Jiaxin Zhang, Lalla Mouatadid, and Kamalika Das.
\newblock {SPUQ}: Perturbation-based uncertainty quantification for large language models.
\newblock In \emph{Proceedings of EACL}, 2024.

\bibitem[Izacard et~al.(2022)]{izacard2022contriever}
Gautier Izacard, Mathilde Caron, Lucas Hosseini, Sebastian Riedel, Piotr Bojanowski, Armand Joulin, and Edouard Grave.
\newblock Unsupervised dense information retrieval with contrastive learning.
\newblock \emph{Transactions on Machine Learning Research}, 2022.

\bibitem[Joshi et~al.(2017)]{joshi2017triviaqa}
Mandar Joshi, Eunsol Choi, Daniel~S. Weld, and Luke Zettlemoyer.
\newblock {TriviaQA}: A large scale distantly supervised challenge dataset for reading comprehension.
\newblock In \emph{Proceedings of ACL}, 2017.

\bibitem[Kadavath et~al.(2022)]{kadavath2022language}
Saurav Kadavath, Tom Conerly, Amanda Askell, Tom Henighan, Dawn Drain, Ethan Perez, Nicholas Schiefer, Zac Hatfield-Dodds, Nova DasSarma, Eli Tyre, et~al.
\newblock Language models (mostly) know what they know.
\newblock \emph{arXiv preprint arXiv:2207.05221}, 2022.

\bibitem[Kossen et~al.(2024)]{kossen2024semantic}
Jannik Kossen, Jiatong Han, Muhammed Razzak, Lisa Schut, Shreshth Malik, and Yarin Gal.
\newblock Semantic entropy probes: Robust and cheap hallucination detection in {LLMs}.
\newblock In \emph{NeurIPS}, 2024.

\bibitem[Kuhn et~al.(2023)]{kuhn2023semantic}
Lorenz Kuhn, Yarin Gal, and Sebastian Farquhar.
\newblock Semantic uncertainty: Linguistic invariances for uncertainty estimation in natural language generation.
\newblock In \emph{ICLR}, 2023.

\bibitem[Kwiatkowski et~al.(2019)]{kwiatkowski2019natural}
Tom Kwiatkowski, Jennimaria Palomaki, Olivia Redfield, Michael Collins, Ankur Parikh, Chris Alberti, Danqi Chen, Eunsol Choi, and Kelvin Guu.
\newblock Natural questions: A benchmark for question answering research.
\newblock \emph{Transactions of the Association for Computational Linguistics}, 2019.

\bibitem[Lewis et~al.(2020)]{lewis2020rag}
Patrick Lewis, Ethan Perez, Aleksandra Piktus, Fabio Petroni, Vladimir Karpukhin, Naman Goyal, Heinrich K{\"u}ttler, Mike Lewis, Wen-tau Yih, Tim Rockt{\"a}schel, Sebastian Riedel, and Douwe Kiela.
\newblock Retrieval-augmented generation for knowledge-intensive {NLP} tasks.
\newblock In \emph{NeurIPS}, 2020.

\bibitem[Mallen et~al.(2023)]{mallen2023when}
Alex Mallen, Akari Asai, Victor Zhong, Rajarshi Das, Daniel Khashabi, and Hannaneh Hajishirzi.
\newblock When not to trust language models: Investigating effectiveness of parametric and non-parametric memories.
\newblock In \emph{Proceedings of ACL}, 2023.

\bibitem[Chen et~al.(2024)]{xia2024controlling}
Lu Chen, Ruqing Zhang, Jiafeng Guo, Yixing Fan, and Xueqi Cheng.
\newblock Controlling risk of retrieval-augmented generation: A counterfactual prompting framework.
\newblock \emph{arXiv preprint arXiv:2409.16146}, 2024.

\bibitem[Shi et~al.(2024)]{shi-etal-2024-trusting}
Weijia Shi, Xiaochuang Han, Mike Lewis, Yulia Tsvetkov, Luke Zettlemoyer, and Wen-tau Yih.
\newblock Trusting your evidence: Hallucinate less with context-aware decoding.
\newblock In \emph{Proceedings of NAACL}, pages 783--791, 2024.

\bibitem[Soudani et~al.(2025)]{soudani2025why}
Heydar Soudani, Evangelos Kanoulas, and Faegheh Hasibi.
\newblock Why uncertainty estimation methods fall short in {RAG}: An axiomatic analysis.
\newblock In \emph{Findings of ACL}, 2025.

\bibitem[Wang et~al.(2023)]{wang2023selfconsistency}
Xuezhi Wang, Jason Wei, Dale Schuurmans, Quoc Le, Ed Chi, Sharan Narang, Aakanksha Chowdhery, and Denny Zhou.
\newblock Self-consistency improves chain of thought reasoning in language models.
\newblock In \emph{ICLR}, 2023.

\bibitem[Yuan et~al.(2026)]{yuan2026crux}
Mingruo Yuan, Shuyi Zhang, and Ben Kao.
\newblock A context-aware dual-metric framework for confidence estimation in large language models.
\newblock \emph{arXiv preprint arXiv:2508.00600}, 2026.

\end{thebibliography}

\newpage
\appendix
\section{Additional Per-Model Results}
\label{app:per_model}

\begin{table}[h]
\caption{AUROC for all \ours{} ablation variants by model (single-seed representative results).}
\label{tab:full_ablation}
\centering
\begin{adjustbox}{max width=\textwidth}
\begin{tabular}{l ccc ccc ccc}
\toprule
& \multicolumn{3}{c}{Llama 3.1 8B} & \multicolumn{3}{c}{Mistral 7B} & \multicolumn{3}{c}{Phi-3 Mini} \\
\cmidrule(lr){2-4} \cmidrule(lr){5-7} \cmidrule(lr){8-10}
Variant & NQ & TQA & PQA & NQ & TQA & PQA & NQ & TQA & PQA \\
\midrule
RS only & 0.595 & 0.587 & 0.700 & 0.617 & 0.816 & 0.706 & 0.550 & 0.519 & 0.667 \\
CD only & 0.502 & 0.505 & 0.578 & 0.495 & 0.632 & 0.521 & 0.526 & 0.616 & 0.508 \\
PA only & 0.481 & 0.503 & 0.558 & 0.476 & 0.511 & 0.607 & 0.474 & 0.611 & 0.680 \\
\midrule
RS+CD & 0.591 & 0.587 & 0.700 & 0.615 & 0.804 & 0.678 & 0.550 & 0.602 & 0.668 \\
RS+PA & 0.592 & 0.579 & 0.697 & 0.595 & 0.791 & 0.758 & 0.479 & 0.625 & 0.644 \\
CD+PA & 0.502 & 0.505 & 0.558 & 0.542 & 0.686 & 0.690 & 0.474 & 0.631 & 0.585 \\
Full & 0.593 & 0.581 & 0.694 & 0.583 & 0.797 & 0.751 & 0.536 & 0.632 & 0.642 \\
\bottomrule
\end{tabular}
\end{adjustbox}
\end{table}

\section{Component Distribution Statistics}
\label{app:components}

\begin{table}[h]
\caption{Component distribution statistics for correct vs.\ incorrect answers (Mistral, TriviaQA).}
\label{tab:component_stats}
\centering
\begin{tabular}{l cc cc}
\toprule
& \multicolumn{2}{c}{Correct} & \multicolumn{2}{c}{Incorrect} \\
\cmidrule(lr){2-3} \cmidrule(lr){4-5}
Component & Mean & Std & Mean & Std \\
\midrule
RS & 0.642 & 0.200 & 0.642 & 0.200 \\
CD & 0.578 & 0.264 & 0.456 & 0.237 \\
PA & 0.608 & 0.456 & 0.724 & 0.409 \\
\bottomrule
\end{tabular}
\end{table}

\end{document}
